\documentclass[12pt]{article}

\usepackage[a4paper, top=2cm, bottom=2cm, left=2.5cm, right=2.5cm]{geometry}
\usepackage{amsmath, amssymb}
\usepackage{tikz}
\usetikzlibrary{arrows.meta,positioning,fit,calc}
\usepackage{multirow}
\usepackage{array}
\usepackage{siunitx}
\usepackage{enumitem}
\usepackage{float}
\usepackage{fontspec}
\usetikzlibrary{decorations.pathmorphing}
\usepackage{pdflscape}
\usepackage{tabularx}
\usepackage{booktabs}
\usepackage{graphicx}
\usepackage{listings}
\usepackage[table]{xcolor}
\usepackage{hyperref}
\usepackage{bookmark}
\usepackage{pdfpages}

% ---------------- FONTS ----------------
\setmainfont[
    UprightFont = Handlee-Regular.ttf,
    AutoFakeBold = 2.5
]{Handlee}

\newfontfamily\headingfont{PT.ttf}
\setmonofont{DejaVu Sans Mono}

% ---------------- HANDWRITTEN MATH ----------------
\usepackage{mathastext}
\Mathastext
\everymath{\displaystyle}

% ---------------- GENERAL ----------------
\setlength{\parindent}{0pt}
\pagenumbering{gobble}
\linespread{1.2}

\newcommand{\mychapter}[2]{%
    \phantomsection
    \hypertarget{#2}{}%
    \bookmark[level=0,dest=#2]{#1}%
    \begin{center}
        {\headingfont\Huge #1}
        \par
        \vspace{4pt}
        \hrule
    \end{center}
}

\newcommand{\mysection}[3]{%
    \phantomsection
    \hypertarget{#3}{}%
    \bookmark[level=1,dest=#3]{#1\space #2}%
    \newpage
    \noindent{\headingfont\Large #1\space #2\par}
    \vspace{4pt}
    \hrule
    \vspace{15pt}
}

\newcommand{\mysubsection}[2]{%
    \vspace{5px}
    \noindent{\headingfont\large #1\space #2\par}
    \vspace{4pt}
    \hrule
    \vspace{15pt}
}

\begin{document}

\mychapter{Chapter 4 : Energy Analysis of Closed System}{6f4y5kdh}

\begin{center}
\renewcommand{\arraystretch}{1.25}
\setlength{\tabcolsep}{8pt}

\begin{tabularx}{\textwidth}{|>{\centering\arraybackslash}p{2.0cm}|X|}
\hline
\textbf{Section} & \textbf{Core Revision Focus} \\
\hline

4.1 & Examines how energy is transferred through the expansion and compression of system boundaries. \\
\hline

4.2 & Establishes how conservation of energy governs processes in closed thermodynamic systems. \\
\hline

4.3 & Introduces specific heats as measures of how thermodynamic energy properties vary with temperature. \\
\hline

4.4 & Develops the temperature-based energy-property relations used to analyze ideal gases. \\
\hline

4.5 & Applies energy-property relations to substances whose volume changes are negligible. \\
\hline


\end{tabularx}
\end{center}

\mysection{4.1}{Moving Boundary Work}{4.1}

\mysubsection{}{Core Relations}

\textbf{Quasi-equilibrium process:} system remains nearly in equilibrium at every instant; a continuous state path and well-defined $P$--$V$ relation exist.

For a piston--cylinder:
\begin{equation*}
\delta W_b=F\,ds=P A\,ds=P\,dV
\end{equation*}

\begin{equation*}
W_b=\int_1^2P\,dV
\end{equation*}

\begin{itemize}
    \item $P$ = absolute pressure.
    \item Expansion: $dV>0\Rightarrow W_b>0$ \; (work by system).
    \item Compression: $dV<0\Rightarrow W_b<0$ \; (work on system).
    \item Boundary work is a \textbf{path function}.
\end{itemize}

For nonquasi-equilibrium motion:
\begin{equation*}
W_b=\int_1^2P_i\,dV
\end{equation*}
where $P_i$ is the pressure acting on the inner face of the piston.

\mysubsection{}{Graphical Interpretation}

For a quasi-equilibrium process:
\begin{equation*}
dA=P\,dV
\end{equation*}

\begin{equation*}
A=\int_1^2P\,dV=W_b
\end{equation*}

\begin{itemize}
    \item Area under a $P$--$V$ process curve $\rightarrow$ magnitude of boundary work.
    \item Area on a $P$--$v$ diagram $\rightarrow$ boundary work per unit mass.
    \item For a cycle, enclosed $P$--$V$ area $\rightarrow$ magnitude of net work.
\end{itemize}

\begin{equation*}
W_{\mathrm{net}}
=
W_{\mathrm{expansion}}-W_{\mathrm{compression}}
\end{equation*}

\mysubsection{}{Experimental and Constant-Volume Cases}

If $P=f(V)$ is unavailable analytically but experimental $P$--$V$ data are known:
\begin{enumerate}
    \item plot the data,
    \item construct the process curve,
    \item determine the area graphically.
\end{enumerate}

Constant volume:
\begin{equation*}
dV=0
\end{equation*}
\begin{equation*}
W_b=0
\end{equation*}

\mysubsection{}{Constant-Pressure Process}

\begin{equation*}
W_b=P_0(V_2-V_1)
\end{equation*}

Using $V=mv$:
\begin{equation*}
W_b=mP_0(v_2-v_1)
\end{equation*}

\mysubsection{}{Isothermal Ideal-Gas Process}

For constant $T$:
\begin{equation*}
PV=mRT_0=C
\end{equation*}

\begin{equation*}
P=\frac{C}{V}
\end{equation*}

Hence,
\begin{equation*}
W_b=C\ln\left(\frac{V_2}{V_1}\right)
\end{equation*}

\begin{equation*}
W_b=P_1V_1\ln\left(\frac{V_2}{V_1}\right)
\end{equation*}

For an ideal gas:
\begin{equation*}
P_1V_1=P_2V_2=mRT_0
\end{equation*}

\begin{equation*}
\frac{V_2}{V_1}=\frac{P_1}{P_2}
\end{equation*}

Compression $\Rightarrow V_2<V_1\Rightarrow W_b<0$.

\mysubsection{}{Polytropic Process}

\begin{equation*}
PV^n=C
\end{equation*}

For $n\neq1$:
\begin{equation*}
W_b=
\frac{P_2V_2-P_1V_1}{1-n}
\end{equation*}

For an ideal gas:
\begin{equation*}
W_b=
\frac{mR(T_2-T_1)}{1-n}
\end{equation*}

\mysubsection{}{Expansion Against a Linear Spring}

Spring force:
\begin{equation*}
F=kx
\end{equation*}

Spring pressure:
\begin{equation*}
P_{\mathrm{spring}}=\frac{F}{A}
\end{equation*}

Since spring force varies linearly with displacement, the $P$--$V$ path is a straight line.

Spring work:
\begin{equation*}
W_{\mathrm{spring}}
=
\frac{1}{2}k(x_2^2-x_1^2)
\end{equation*}

Total work:
\begin{equation*}
W_{\mathrm{total}}
=
P_1(V_2-V_1)
+
\frac{1}{2}k(x_2^2-x_1^2)
\end{equation*}

\mysubsection{}{Energy Destination of Boundary Work}

Boundary work transfers energy to other mechanisms, e.g.
\begin{equation*}
W_b=
W_{\mathrm{friction}}
+
W_{\mathrm{atm}}
+
W_{\mathrm{crank}}
\end{equation*}

\mysection{4.2}{Energy Balance for Closed Systems}{4.2}

\mysubsection{}{General Energy Balance}

Conservation of energy:
\begin{equation*}
E_{\mathrm{in}}-E_{\mathrm{out}}
=
\Delta E_{\mathrm{system}}
\end{equation*}

Total system energy may include internal, kinetic, potential, and other forms.

Rate form:
\begin{equation*}
\dot E_{\mathrm{in}}-\dot E_{\mathrm{out}}
=
\frac{dE_{\mathrm{system}}}{dt}
\end{equation*}

For constant rates:
\begin{equation*}
Q=\dot Q\,\Delta t,
\qquad
W=\dot W\,\Delta t,
\qquad
\Delta E=
\frac{dE}{dt}\Delta t
\end{equation*}

Per unit mass:
\begin{equation*}
e_{\mathrm{in}}-e_{\mathrm{out}}
=
\Delta e_{\mathrm{system}}
\end{equation*}

\mysubsection{}{Cycle and First-Law Forms}

For a cycle:
\begin{equation*}
\Delta E_{\mathrm{system}}=0
\end{equation*}

\begin{equation*}
E_{\mathrm{in}}=E_{\mathrm{out}}
\end{equation*}

For a closed-system cycle:
\begin{equation*}
W_{\mathrm{net,out}}
=
Q_{\mathrm{net,in}}
\end{equation*}

\begin{equation*}
\dot W_{\mathrm{net,out}}
=
\dot Q_{\mathrm{net,in}}
\end{equation*}

\textbf{Classical sign convention:}
\begin{itemize}
    \item $Q>0$: heat transfer into system.
    \item $W>0$: work done by system.
\end{itemize}

Thus:
\begin{equation*}
Q-W=\Delta E
\end{equation*}

where
\begin{equation*}
Q=Q_{\mathrm{in}}-Q_{\mathrm{out}}
\end{equation*}
and
\begin{equation*}
W=W_{\mathrm{out}}-W_{\mathrm{in}}
\end{equation*}

Negative $Q$ or $W$ means the actual direction is opposite to the assumed positive direction.

\mysubsection{}{Stationary Closed System}

If kinetic and potential energy changes are negligible:
\begin{equation*}
Q-W=\Delta U
\end{equation*}

Per unit mass:
\begin{equation*}
q-w=\Delta e
\end{equation*}

Differential form:
\begin{equation*}
\delta q-\delta w=de
\end{equation*}

Heat and work are both \textbf{energy-transfer mechanisms}, not properties of the system.

\mysubsection{}{Unrestrained Expansion}

Expansion into a vacuum:
\begin{equation*}
P_{\mathrm{ext}}=0
\end{equation*}

Therefore:
\begin{equation*}
W_b=\int_1^2P_{\mathrm{ext}}\,dV=0
\end{equation*}

If no other work is present:
\begin{equation*}
W=0
\end{equation*}

For a stationary closed system:
\begin{equation*}
Q=\Delta U
\end{equation*}

\textbf{Important:} work depends on the chosen system boundary; including the evacuated space can make the boundary fixed.

\mysection{4.3}{Specific Heats}{4.3}

\mysubsection{}{Definitions}

Specific heat = energy required to raise the temperature of unit mass by one degree under a specified condition.

\begin{equation*}
c_v=
\left(\frac{\partial u}{\partial T}\right)_v
\end{equation*}

\begin{equation*}
c_p=
\left(\frac{\partial h}{\partial T}\right)_p
\end{equation*}

Hence:
\begin{itemize}
    \item $c_v$ describes the variation of internal energy with $T$ at constant $v$.
    \item $c_p$ describes the variation of enthalpy with $T$ at constant $P$.
    \item Generally, $c_p>c_v$.
\end{itemize}

These are \textbf{property relations}; they are not restricted to processes actually occurring at constant $v$ or constant $P$.

\mysubsection{}{State Dependence and Units}

In general, specific heats depend on the state of the substance.

Common units:
\begin{equation*}
\mathrm{\frac{kJ}{kg\cdot{}^\circ C}}
\qquad\text{or}\qquad
\mathrm{\frac{kJ}{kg\cdot K}}
\end{equation*}

Since temperature intervals have equal magnitude:
\begin{equation*}
\Delta T(^\circ\mathrm C)=\Delta T(\mathrm K)
\end{equation*}

Molar specific heats:
\begin{equation*}
\bar c_v,\qquad \bar c_p
\end{equation*}

\mysection{4.4}{Internal Energy, Enthalpy, and Specific Heats of Ideal Gases}{4.4}

\mysubsection{}{Ideal-Gas Relations}

\begin{equation*}
Pv=RT
\end{equation*}

For an ideal gas:
\begin{equation*}
u=u(T)
\end{equation*}

Enthalpy definition:
\begin{equation*}
h=u+Pv
\end{equation*}

Since $Pv=RT$:
\begin{equation*}
h=u+RT
\end{equation*}

Therefore:
\begin{equation*}
h=h(T)
\end{equation*}

\textbf{Key result:} for an ideal gas, $u$, $h$, $c_v$, and $c_p$ depend only on temperature.

\mysubsection{}{Ideal-Gas Specific Heats}

\begin{equation*}
c_v=c_v(T)
\end{equation*}

\begin{equation*}
c_p=c_p(T)
\end{equation*}

Differential relations:
\begin{equation*}
du=c_v(T)\,dT
\end{equation*}

\begin{equation*}
dh=c_p(T)\,dT
\end{equation*}

Changes between states:
\begin{equation*}
\Delta u=
\int_{T_1}^{T_2}c_v(T)\,dT
\end{equation*}

\begin{equation*}
\Delta h=
\int_{T_1}^{T_2}c_p(T)\,dT
\end{equation*}

Thus, $\Delta u$ and $\Delta h$ depend only on $T_1,T_2$, not on the process path.

\mysubsection{}{Average Specific Heat}

For moderate/small temperature ranges:
\begin{equation*}
\Delta u
\approx
c_{v,\mathrm{avg}}(T_2-T_1)
\end{equation*}

\begin{equation*}
\Delta h
\approx
c_{p,\mathrm{avg}}(T_2-T_1)
\end{equation*}

Useful temperature estimate:
\begin{equation*}
T_{\mathrm{avg}}
=
\frac{T_1+T_2}{2}
\end{equation*}

or
\begin{equation*}
c_{\mathrm{avg}}
=
\frac{c_1+c_2}{2}
\end{equation*}

\textbf{Important:}
\begin{equation*}
\Delta u=c_{v,\mathrm{avg}}\Delta T
\end{equation*}
and
\begin{equation*}
\Delta h=c_{p,\mathrm{avg}}\Delta T
\end{equation*}
are applicable to \textbf{any ideal-gas process}; the subscripts do not specify the process path.

\mysubsection{}{Monatomic Ideal Gases}

Specific heats are essentially constant over a wide range:
\begin{equation*}
\Delta u=c_v(T_2-T_1)
\end{equation*}

\begin{equation*}
\Delta h=c_p(T_2-T_1)
\end{equation*}

\mysubsection{}{Methods for Determining $\Delta u$ and $\Delta h$}

\begin{enumerate}
    \item \textbf{Property tables}
    \begin{equation*}
    \Delta u=u_2-u_1,
    \qquad
    \Delta h=h_2-h_1
    \end{equation*}

    \item \textbf{Temperature-dependent specific heats}
    \begin{equation*}
    \Delta u=\int_{T_1}^{T_2}c_v(T)\,dT,
    \qquad
    \Delta h=\int_{T_1}^{T_2}c_p(T)\,dT
    \end{equation*}

    \item \textbf{Average specific heats}
    \begin{equation*}
    \Delta u\approx c_{v,\mathrm{avg}}\Delta T,
    \qquad
    \Delta h\approx c_{p,\mathrm{avg}}\Delta T
    \end{equation*}
\end{enumerate}

\mysubsection{}{Reference State}

Reference values of $u$ and $h$ are arbitrary; only differences matter:
\begin{equation*}
\Delta u=u_2-u_1
\end{equation*}

\begin{equation*}
\Delta h=h_2-h_1
\end{equation*}

\mysubsection{}{Specific-Heat Relations}

From
\begin{equation*}
h=u+RT
\end{equation*}

\begin{equation*}
dh=du+R\,dT
\end{equation*}

Therefore:
\begin{equation*}
c_p=c_v+R
\end{equation*}

\begin{equation*}
c_v=c_p-R
\end{equation*}

On a molar basis:
\begin{equation*}
\bar c_p=\bar c_v+R_u
\end{equation*}

Specific heat ratio:
\begin{equation*}
k=\frac{c_p}{c_v}
\end{equation*}

$k$ varies mildly with temperature and may often be treated as constant over a limited temperature range.

\mysection{4.5}{Internal Energy, Enthalpy, and Specific Heats of Solids and Liquids}{4.5}

\mysubsection{}{Incompressible Substances}

For an incompressible substance:
\begin{equation*}
v=\text{constant}
\end{equation*}

or
\begin{equation*}
\rho=\text{constant}
\end{equation*}

Solids and liquids are commonly approximated as incompressible.

\mysubsection{}{Specific Heat and Internal Energy}

For incompressible substances:
\begin{equation*}
c_p=c_v=c
\end{equation*}

and
\begin{equation*}
c=c(T)
\end{equation*}

Internal-energy change:
\begin{equation*}
du=c(T)\,dT
\end{equation*}

\begin{equation*}
\Delta u=
\int_{T_1}^{T_2}c(T)\,dT
\end{equation*}

For approximately constant $c$:
\begin{equation*}
\Delta u
\approx
c_{\mathrm{avg}}(T_2-T_1)
\end{equation*}

\mysubsection{}{Enthalpy of Incompressible Substances}

\begin{equation*}
h=u+Pv
\end{equation*}

Since
\begin{equation*}
dv=0
\end{equation*}

then
\begin{equation*}
dh=du+v\,dP
\end{equation*}

Using $du=c(T)\,dT$:
\begin{equation*}
dh=c(T)\,dT+v\,dP
\end{equation*}

Thus:
\begin{equation*}
\Delta h=\Delta u+v\Delta P
\end{equation*}

Average-$c$ form:
\begin{equation*}
\Delta h
\approx
c_{\mathrm{avg}}\Delta T+v\Delta P
\end{equation*}

\begin{equation*}
\Delta h
\approx
c_{\mathrm{avg}}(T_2-T_1)+v(P_2-P_1)
\end{equation*}

\mysubsection{}{Special Cases}

\textbf{Solids:}
\begin{equation*}
v\Delta P\approx0
\end{equation*}

Therefore:
\begin{equation*}
\Delta h\approx\Delta u
\end{equation*}

\begin{equation*}
\Delta h
\approx
c_{\mathrm{avg}}(T_2-T_1)
\end{equation*}

\textbf{Liquid at constant pressure:}
\begin{equation*}
\Delta P=0
\end{equation*}

\begin{equation*}
\Delta h=\Delta u
\end{equation*}

\begin{equation*}
\Delta h
\approx
c_{\mathrm{avg}}(T_2-T_1)
\end{equation*}

\textbf{Liquid at constant temperature:}
\begin{equation*}
\Delta T=0
\end{equation*}

\begin{equation*}
\Delta u=0
\end{equation*}

\begin{equation*}
\Delta h=v\Delta P
\end{equation*}

\mysubsection{}{Compressed Liquid Enthalpy}

Approximation from saturated-liquid properties at the same temperature:
\begin{equation*}
h_{P,T}
\approx
h_f(T)
+
v_f(T)
\left[P-P_{\mathrm{sat}}(T)\right]
\end{equation*}

Simpler approximation:
\begin{equation*}
h_{P,T}\approx h_f(T)
\end{equation*}

The pressure correction may be neglected when sufficiently small; at high $T$ and $P$, the simpler approximation can sometimes be preferable.

\mysubsection{}{Energy Balance for Multiple Incompressible Substances}

For several substances:
\begin{equation*}
\Delta U_{\mathrm{sys}}
=
\sum_i\Delta U_i
\end{equation*}

For an isolated system:
\begin{equation*}
\Delta U_{\mathrm{sys}}=0
\end{equation*}

For each substance:
\begin{equation*}
\Delta U_i=m_i\Delta u_i
\end{equation*}

With constant specific heats:
\begin{equation*}
\Delta U_i
\approx
m_i c_i(T_2-T_{1,i})
\end{equation*}

Hence:
\begin{equation*}
\boxed{
\sum_i m_i c_i(T_2-T_{1,i})=0
}
\end{equation*}

For thermal equilibrium:
\begin{equation*}
T_1=T_2=\cdots=T_{\mathrm{eq}}
\end{equation*}

Energy lost by hotter substances = energy gained by colder substances.

\mysubsection{}{Heat Transfer for Incompressible Substances}

For constant $c$:
\begin{equation*}
\Delta U=mc(T_2-T_1)
\end{equation*}

Cooling:
\begin{equation*}
Q_{\mathrm{out}}
=
mc(T_1-T_2)
\qquad (T_1>T_2)
\end{equation*}

Heating:
\begin{equation*}
Q_{\mathrm{in}}
=
mc(T_2-T_1)
\qquad (T_2>T_1)
\end{equation*}


\end{document}